\begin{table}[t]
\centering
\caption{Paired test-set evidence for the refine-network gain over the pseudo-label.}
\label{tab:refine-statistics}
\begin{threeparttable}
\begin{tabular}{@{}llc@{}}
\toprule
Evidence type & Measurement & Value \\
\midrule
\multicolumn{3}{@{}l}{\textit{Primary paired comparison on 31 held-out scans}} \\
Input & Pseudo-label DSC & $0.9278 \pm 0.0505$ \\
Output & Refined DSC & $0.9341 \pm 0.0453$ \\
Gain & Mean paired $\Delta$DSC & $+0.0063 \pm 0.0076$ \\
Gain & Median paired $\Delta$DSC & $+0.0076$ \\
Direction & Improved / worsened scans & 25 / 6 \\
\addlinespace[2pt]
\multicolumn{3}{@{}l}{\textit{Significance of the paired gain}} \\
Parametric test & Paired $t$-test, one-sided & $p=3.46\times 10^{-5}$ \\
Rank test & Wilcoxon signed-rank, one-sided & $p=4.33\times 10^{-5}$ \\
Direction test & Sign test, one-sided & $p=4.39\times 10^{-4}$ \\
Bootstrap & 95\% CI for mean $\Delta$DSC & $[0.0037,\ 0.0090]$ \\
Standardized effect & Cohen's $d_z$ & $0.829$ \\
\bottomrule
\end{tabular}
\begin{tablenotes}[flushleft]\footnotesize
\item DSC: Dice similarity coefficient. The table is intentionally limited to paired test-set evidence because the refine network is evaluated by whether it improves its own input pseudo-label on the same scans. The checkpoint (epoch 60) and probability threshold (0.70) were selected on the validation split before test evaluation; test labels were used only for the metrics and paired tests reported here.
\end{tablenotes}
\end{threeparttable}
\end{table}
